% ===== §19 The Aesthetic Everyday and Aestheticized Consumption =====
\section{The Aesthetic Everyday and Aestheticized Consumption}
\label{p22:sec:fetishism}

\noindent
The account of aesthetic education developed here has a counterfeit, and the counterfeit is close enough to be mistaken for it, which makes the distinction between them urgent rather than merely tidy. The counterfeit is the aestheticization of the everyday that the contemporary market promotes under the names of refined living, curated domesticity, the beautiful life, in which the rhythms and objects and gestures of daily life are made objects of display and consumption. This looks like what the paper recommends, the bringing of aesthetic attention to the everyday, and it is in fact its inversion, the wheel dressed as the spiral. This section draws the distinction, using the criterion the series has carried throughout, the difference between the cycle that gains and the cycle that turns in place, and it locates the counterfeit within the three-layer account of value the series has developed, showing aestheticized consumption to be the everyday's holonomy captured in the imaginary register.

\subsection{The counterfeit and its appeal}
\label{p22:sub:fetish-counterfeit}

The aestheticization of the everyday is a powerful counterfeit because it borrows the true coin's appearance. It attends to the daily, as the paper urges; it cares about the placement of things, the quality of the meal, the texture of the morning, as field-building does; it presents itself as the refinement of ordinary life into something fine and worth living. Its appeal is real and not simply a deception, because the impulse it answers is a true one, the impulse to bring care and attention to the everyday rather than letting it pass unregarded. What makes it a counterfeit is not that the attention is false but that the attention is turned to the wrong object and serves the wrong cycle. The aestheticized everyday makes the daily an object of display and consumption, a performance of refinement to be seen and a pleasure to be had, and in doing so it removes the daily from the relational field where field-building places it and installs it in a circuit of presentation and acquisition. The care is real; its direction is inverted; and the inversion is exactly the substitution of a wheel for a spiral, a cycle of consumption that turns in place for a cycle of relation that gains.

\subsection{The criterion: which cycle does the attention serve}
\label{p22:sub:fetish-criterion}

The distinction between the aesthetic everyday and aestheticized consumption is drawn by the criterion the series has used for cycles throughout, applied now to the attention the everyday receives. The question is which cycle the attention serves, the generative cycle of the relational field or the consumptive cycle of display and acquisition. Attention to the everyday that serves the relational field is field-building: it perceives the field's state, generates the appropriate act, sustains the accumulation, and the beauty it produces is the coupling event that the field gains. Attention to the everyday that serves the consumptive cycle is aestheticization: it perceives the everyday as an object to be made fine and displayed, generates a performance of refinement, sustains a circuit of acquisition in which each acquired refinement must be succeeded by another, and the beauty it produces is a possession to be had and shown. The two attentions look alike, both careful, both directed at the daily, both producing something called beauty, and they are told apart only by which cycle they serve, the generative one that gains or the consumptive one that turns.

\claim{The aesthetic everyday and aestheticized consumption are distinguished by which cycle the attention serves. Attention that serves the relational field is field-building, perceiving the field, generating the appropriate act, sustaining the accumulation, and the beauty it produces is the coupling event the field gains. Attention that serves the cycle of display and acquisition is aestheticization, making the everyday a possession to be refined and shown, and the beauty it produces is a thing to be had. The two are alike in their care for the daily and are told apart only by their cycle, the generative one that gains against the consumptive one that turns in place. Aestheticized consumption is the everyday made into a wheel.}

\begin{casebox}[Two carefully laid tables]
Two tables are laid with equal care, and from a photograph they cannot be told apart. The first is laid for the other who will sit at it, its care a perception of her, this plate here because she reaches with that hand, this dimness because the day pressed on her, the laying an act that meets her state and builds the field of the evening they will share; the beauty is in the coupling the table prepares, and it will pass when the evening passes, having done its work in the between. The second is laid to be a beautiful table, its care a performance of refinement to be photographed, admired, matched against other tables, the placement governed by an image of the beautiful table and not by the one who will sit at it, who is in a sense incidental to a table that would be equally beautiful with anyone or no one before it. The first table serves the field; the second serves the image. The photograph cannot tell them apart because the difference is not in the table but in the cycle the laying serves, and the same exquisite care builds a field in the first case and feeds a circuit of display in the second. The second table will need to be succeeded by a finer one, because the image it serves is never satisfied; the first needs no successor, because the field it built has gained and will be returned to tomorrow, the same table laid again for the same person, thickening.
\end{casebox}

\subsection{The counterfeit in the register of the imaginary}
\label{p22:sub:fetish-imaginary}

The counterfeit can be located precisely within the three-layer account of value the series developed, and the location explains both its appeal and its exhaustion. That account distinguished the imaginary value of the idolised possession, which exhausts because it must be clung to and never satisfies; the symbolic value of the exchangeable, which is the most easily stolen; and the real value that gathers around the central absence and is the only locus of genuine positive accumulation. Aestheticized consumption is the everyday captured in the imaginary register. It makes the daily an idolised possession, a refinement to be had and clung to, and like all imaginary value it exhausts, which is why it demands an endless succession of new refinements, each consumed and found wanting, each requiring a further acquisition to sustain the feeling it briefly gave. This is the consumptive wheel in its precise form: the imaginary value of the aestheticized everyday cannot accumulate, because imaginary value never accumulates, and so it must be endlessly renewed by fresh consumption, the new curtain, the new ritual, the new performance of the beautiful life, none of which gains because none of them touches the real value that gathers only in the relational field.

The counterfeit is hard to detect from within, and saying why gives the section its sharpest practical point, because the difficulty is structural and not a matter of insufficient honesty. The aestheticized everyday is hard to tell from the aesthetic everyday from within because the imaginary register presents itself as fullness, because the idolised possession feels, in the moment of its acquisition, like the real thing it counterfeits; the new and beautiful object does give a real pleasure, and the performance of the refined life is genuinely pleasing to perform, so that nothing in the immediate feeling distinguishes the consumption from the relation. The difference shows only over time, in whether the pleasure accumulates or must be renewed, and this is the diagnostic sign the section can offer: if the beautiful thing, once had, is enough, if the field it belongs to is thereby richer and the next morning begins higher, the attention served the relation; if the beautiful thing, once had, must be succeeded by another to sustain the feeling, if the satisfaction drains and demands replacement, the attention served the consumptive wheel. The sign is not in the object or the act, which are identical across the two cases, but in the temporal signature, the accumulation or the exhaustion, the gaining or the endless renewal. A person uncertain whether their care for the everyday is field-building or aestheticization can ask not how it feels, which does not distinguish them, but whether it accumulates, whether last month's beautiful evening enriched this month's or merely set the bar for a beauty this month must exceed. The exhaustion that demands ever-more is the signature of the imaginary; the quiet sufficiency that returns to the same and finds it thickened is the signature of the real. This is the same temporal criterion the whole paper has used, the distinction of the spiral from the wheel, applied now as a practical test by which a person may detect, in their own care for the everyday, which cycle it serves.

The aesthetic everyday this paper recommends operates in the other register entirely. Its beauty is the real value that gathers around the coupling, the accumulation that the relational field gains, and because it is real value it accumulates, returns to its root raised, requires no endless succession of new objects because it is not consumed in the having but generated in the relating. The difference between the two is therefore not a difference of taste or degree, the modest everyday against the lavish one, but a difference of register, the real against the imaginary, the accumulating against the exhausting. This is why the aesthetic everyday is available to those with nothing to acquire and the aestheticized everyday is not: field-building requires no new objects, only the relating that generates real value in whatever the field already contains, whereas aestheticized consumption requires a stream of acquisitions and so is available only to those who can afford the stream. The counterfeit is not merely the false version of the true; it is the true version's exclusion of exactly those the true version includes, the dispossessed who cannot consume but can relate, and the distinction the section draws is therefore not only metaphysical but, as the next section will press, political.
